%!TEX root = ../main.tex
% GLOSARIO: términos que necesitan definición dentro del documento.
%
% Definir:   \newglossaryentry{clave}{name={Término}, description={Definición}}
% Usar:      \gls{clave}      -> Término (con enlace al glosario)
%            \Gls{clave}      -> Término (con mayúscula inicial)
%            \glspl{clave}    -> plural
% Solo aparecen en el glosario los términos que se usan en el texto.
% Las siglas van en siglas.tex y los símbolos en simbolos.tex.

\newglossaryentry{g:mvp}{
    name={producto mínimo viable},
    description={Versión funcional de un producto con el conjunto mínimo
    de características que entrega valor real a los usuarios y permite
    aprender de su uso. A diferencia de un prototipo, no simula: funciona}
}

\newglossaryentry{prototipo}{
    name={prototipo},
    description={Versión preliminar de un sistema construida para validar
    una idea, un flujo o una interfaz. Puede simular partes del
    comportamiento y no está pensada para operar en producción}
}

\newglossaryentry{stakeholder}{
    name={parte interesada},
    description={Persona, grupo u organización que puede afectar o verse
    afectada por el proyecto: usuarios, clientes, patrocinadores,
    reguladores y el propio equipo}
}

\newglossaryentry{exposicion}{
    name={exposición al riesgo},
    description={Producto de la probabilidad de ocurrencia de un riesgo por
    su impacto. Se usa para ordenar la matriz de riesgos y decidir cuáles
    se tratan primero}
}

\newglossaryentry{g:van}{
    name={valor actual neto},
    description={Suma de los flujos de caja futuros de un proyecto
    descontados a una tasa que refleja su riesgo, menos la inversión
    inicial. Un valor positivo indica que el proyecto crea valor}
}

\newglossaryentry{matching}{
    name={matching},
    description={Proceso por el cual el sistema cruza los perfiles de
    alumnos y tutores para recomendar la combinación más compatible, en
    el caso de Tinku mediante similitud semántica de sus descripciones}
}

\newglossaryentry{marketplace}{
    name={marketplace},
    description={Plataforma digital bilateral que conecta a dos grupos de
    usuarios —en este caso, estudiantes que demandan tutorías y tutores
    que las ofrecen— y facilita la transacción entre ellos}
}

\newglossaryentry{embeddings}{
    name={embeddings semánticos},
    description={Representaciones vectoriales de un texto en un espacio de
    alta dimensionalidad, generadas por modelos de lenguaje, que capturan
    el significado de las oraciones y permiten compararlas por similitud}
}

\newglossaryentry{escrow}{
    name={escrow},
    description={Mecanismo de pago en el que el dinero del comprador se
    retiene en un intermediario y solo se libera al vendedor cuando la
    contraprestación se efectiviza, reduciendo el riesgo de fraude en
    plataformas bilaterales}
}

\newglossaryentry{coldstart}{
    name={arranque en frío},
    description={Problema inicial de los marketplaces bilaterales: la
    plataforma solo es valiosa cuando hay usuarios de ambos lados, pero
    la falta de usuarios iniciales dificulta atraer a los primeros}
}
